\documentclass[11pt, a4paper, titlepage]{article}
\usepackage{amsmath}
\usepackage{a4wide}
\usepackage{url}
\usepackage{multirow}
\usepackage{amsfonts}
\newcommand{\release}{5.0.0}
\newcommand{\JAGS}{\textsf{JAGS}}
\newcommand{\Rtools}{\textsf{Rtools45}}
\newcommand{\BUGS}{\textsf{BUGS}}
\newcommand{\WinBUGS}{\textsf{WinBUGS}}
\newcommand{\R}{\textsf{R}}
\newcommand{\CODA}{\textsf{coda}}

\usepackage{verbatim}

\newcommand{\code}[1]{{\bgroup{\normalfont\ttfamily #1}\egroup}}
\newcommand{\samp}[1]{{`\bgroup\normalfont\texttt{#1}'\egroup}}
\newcommand{\file}[1]{{`\normalfont\textsf{#1}'}}
\let\command=\code
\let\option=\samp

\begin{document}

\title{JAGS Version \release\ installation manual}
\author{Martyn Plummer \and Matt Denwood}
\date{18 December 2025}
\maketitle

\JAGS\ is distributed in binary format for Microsoft Windows, macOS, 
and most Linux distributions.  The following instructions are for
those who wish to build \JAGS\ from source. The manual is divided
into three sections with instructions for Linux/Unix, macOS, and Windows.

\section{Linux and UNIX}

\JAGS\ follows the usual GNU convention of 
\begin{verbatim}
./configure
make
make install
\end{verbatim}
which is described in more detail in the file \texttt{INSTALL} in
the top-level source directory. On some UNIX platforms, you may
be required to use GNU make (gmake) instead of the native make
command. On systems with multiple processors, you may use the option 
\option{-j} to speed up compilation, {\em e.g.} for a quad-core PC you
may use:
\begin{verbatim}
make -j4
\end{verbatim}
If you have the cppunit library installed then you can test the build with
\begin{verbatim}
make check
\end{verbatim}
WARNING. If you already have a copy of the jags library installed on
your system then the test program created by \code{make check} will
link against the {\bf installed} library and not the one in your build
directory. So if the test suite includes a regression test for a bug
that was fixed in the version you are building but a previous version
of JAGS is already installed then the unit tests will fail. Best
practice is to run the tests after \code{make install} so the build
and installed versions are the same.

\subsection{Configure options}

At configure time you also have the option of defining options such
as:
\begin{itemize}
\item The names of the C, C++, and Fortran compilers.  
\item Optimization flags for the compilers.  \JAGS\ is optimized by
  default if the GNU compiler (gcc) is used. If you are using another
  compiler then you may need to explicitly supply optimization flags.
\item Installation directories. \JAGS\ conforms to the GNU standards
  for where files are installed. You can control the installation
  directories in more detail using the flags that are listed when
  you type \command{./configure --help}.
\end{itemize}

\subsubsection{Configuration for a 64-bit build}

By default, JAGS will install all libraries into
\file{/usr/local/lib}.  If you are building a 64-bit version of \JAGS,
this may not be appropriate for your system. On Fedora and other
RPM-based distributions, for example, 64-bit libraries should be
installed in \file{lib64}, and on Solaris, 64-bit libraries are in a
subdirectory of \file{lib} ({\em e.g.} \file{lib/amd64} if you are
using a x86-64 processor), whereas on Debian, and other Linux
distributions that conform to the FHS, the correct installation
directory is \file{lib}.

To ensure that \JAGS\ libraries are installed in the correct
directory, you should supply the \option{--libdir} argument to the
configure script, {\em e.g.}:
\begin{verbatim}
./configure --libdir=/usr/local/lib64
\end{verbatim}

It is important to get the installation directory right when using the
\code{rjags} interface between R and \JAGS, otherwise the
\code{rjags} package will not be able to find the \JAGS\ library.

\subsubsection{Configuration for a private installation}

If you do not have administrative privileges, you may wish to install
\JAGS\ in your home directory. This can be done with the following
configuration options
\begin{verbatim}
export JAGS_HOME=$HOME/jags #or wherever you want it
./configure --prefix=$JAGS_HOME
\end{verbatim}
For more detailed control over the installation directories type
\begin{verbatim}
./configure --help
\end{verbatim}
and read the section ``Fine-tuning of the installation directories.''

With a private installation, you need to modify your PATH environment
variable to include \file{\$JAGS\_HOME/bin}. You may also need to set
\code{LD\_LIBRARY\_PATH} to include \file{\$JAGS\_HOME/lib} (On Linux this
is not necessary as the location of libjags and libjrmath is hard-coded
into the \JAGS\ binary).

\subsection{BLAS and LAPACK}
\label{section:blas:lapack}

BLAS (Basic Linear Algebra System) and LAPACK (Linear Algebra Pack)
are two libraries of routines for linear algebra. They are used by the
multivariate functions and distributions in the \texttt{bugs} module.
Most unix-like operating system vendors supply shared libraries that
provide the BLAS and LAPACK functions, although the libraries may not
literally be called ``blas'' and ``lapack''.  During configuration, a
default list of these libraries will be checked. If \texttt{configure}
cannot find a suitable library, it will stop with an error message.

You may use alternative BLAS and LAPACK libraries using the configure
options \texttt{--with-blas} and \texttt{--with-lapack}
\begin{verbatim}
./configure --with-blas="-lmyblas" --with-lapack="-lmylapack"
\end{verbatim}

If the BLAS and LAPACK libraries are in a directory that is not on the
default linker path, you must set the \code{LDFLAGS} environment variable
to point to this directory at configure time:
\begin{verbatim}
LDFLAGS="-L/path/to/my/libs" ./configure ...
\end{verbatim}

At runtime, if you have linked \JAGS\ against BLAS or LAPACK in
a non-standard location, you must supply this location with the
environment variable \code{LD\_LIBRARY\_PATH}, {\em e.g.}
\begin{verbatim}
LD_LIBRARY_PATH="/path/to/my/libs:${LD_LIBRARY_PATH}"
\end{verbatim}
Since the configure script will attempt to compile and run test
programs linked to BLAS and LAPACK, you may also need to set
\code{LD\_LIBRARY\_PATH} at configure time.

An alternative to setting \code{LD\_LIBRARY\_PATH} is to hard-code the
paths to the blas and lapack libraries at compile time. This is
compiler and platform-specific, but is typically achieved with
\begin{verbatim}
LDFLAGS="-L/path/to/my/libs -R/path/to/my/libs
\end{verbatim}

JAGS can also be linked to static BLAS and LAPACK if they have both
been compiled with the \texttt{-fPIC} flag. You will probably need to
do a custom build of BLAS and LAPACK if you require this. The
configure options for JAGS are then:
\begin{verbatim}
./configure --with-blas="-L/path/to/my/libs -lmyblas -lgfortran -lquadmath" \
            --with-lapack="-L/path/to/my/libs -lmylapack"
\end{verbatim}
Note that with static linking you must add the dependencies of the BLAS
library manually. The LAPACK library will pick up the same dependencies.
Note also that libtool does not like linking directly to archive files.
If you attempt a configuration of the form
\begin{verbatim}
--with-blas="/path/to/my/libs/myblas.a"
\end{verbatim}
then this will pass at configure time but ``make'' will not correctly
build the JAGS modules.


\subsection{GNU/Linux}
\label{section:gnulinux}

GNU/Linux is the development platform for \JAGS, and a variety of
different build options have been explored, including the use of
third-party compilers and linear algebra libraries.

\subsubsection{BLAS and LAPACK}

The {\bf BLAS} and {\bf LAPACK} libraries from Netlib
(\url{http://www.netlib.org}) should be provided as part of your Linux
distribution. If your Linux distribution splits packages into ``user''
and ``developer'' versions, then you must install the developer
package ({\em e.g.}  \texttt{blas-devel} and \texttt{lapack-devel}, or
\texttt{blas-dev} and \texttt{lapack-dev}).

{\bf OpenBLAS} is an optimized BLAS library that may be used as a
replacement for Netlib BLAS.

{\bf BLIS} is a portable framework for instantiating BLAS-like
libraries for dense linear algebra operations. While BLIS has its own
API it also includes a BLAS compatibility layer and can be used as a
drop-in replacement for BLAS. The same team that created BLIS also
created {\bf libflame}, which provides much of the functionality of
LAPACK and can also be used with \JAGS.

OpenBLAS and BLIS may come in serial and parallel versions on your
Linux distribution. We recommend using the serial versions for \JAGS.

\subsubsection{FlexiBLAS}

FlexiBLAS is a front end for BLAS/LAPACK. FlexiBLAS is the default for
current RedHat based Linux distributions (Fedora and RHEL).

At configuration time, \JAGS\ will look for FlexiBLAS first before any
other BLAS implementation.  If \JAGS\ is linked against FlexiBLAS then
different BLAS/LAPACK backends can be chosen at runtime. A list of
available backends can be obtained with the command line interface:
\begin{verbatim}
flexiblas list
\end{verbatim}
You can then set the default BLAS backend from this list, {\em e.g.}
\begin{verbatim}
flexiblas default ATLAS
\end{verbatim}
Or you can temporarily set the backend with the environment variable
\texttt{FLEXIBLAS}, {\em e.g.}
\begin{verbatim}
export FLEXIBLAS=NETLIB
\end{verbatim}
If using an R interface to \texttt{JAGS}, you can use the functions in
the \texttt{flexiblas} package, which is available from CRAN, to set
the backend and to control the number of threads used by multithreaded
backends.

\subsubsection{AMD Optimizing CPU Libraries}
\label{section:aocl:linux}

The AMD Optimizing CPU Libraries (AOCL) include a version of the BLIS
and libflame libraries for AMD processors. To link \JAGS\ with
\texttt{AOCL}, you must supply \texttt{--with-blas=-lblis} and
\texttt{--with-lapack=-lflame}.  Since AOCL is by default installed
into a private location you will also need to set the environment
variables \texttt{LD\_LIBRARY\_PATH} and \texttt{LDFLAGS}
appropriately, {\em e.g.}

\begin{verbatim}
LD_LIBRARY_PATH=${LD_LIBRARY_PATH}:${HOME}/aocl/5.1.0/gcc/lib
LDFLAGS="-L${HOME}/aocl/5.1.0/gcc/lib \
./configure --with-blas="-lblis" --with-lapack="-lflame"
\end{verbatim}

\subsubsection{Intel oneAPI Math Kernel Library}

The Intel oneAPI Math Kernel library (MKL) provides optimized BLAS and
LAPACK routines for Intel processors.

MKL uses a layered model for linking which gives the user fine-grained
control over four different library layers: interface, threading,
computation, and run-time. Some examples of linking to MKL using this
layered model are given below. For further guidance, consult the MKL
Link Line advisor at
\url{https://www.intel.com/content/www/us/en/developer/tools/oneapi/onemkl-link-line-advisor.html}

MKL includes a shell script that sets up the environment variables
necessary to build an application with MKL.
\begin{verbatim}
source /opt/intel/oneapi/setvars.sh
\end{verbatim}

After calling this script, you can link \JAGS\ with a sequential
version of MKL as follows:
\begin{verbatim}
./configure \
--with-blas="-lmkl_intel_lp64 -lmkl_sequential -lmkl_core -lpthread -lm -ldl"
\end{verbatim}
Note that \code{libpthread} is still required, even when linking
to sequential MKL.

MKL with OpenMP threads may be used with:
\begin{verbatim}
./configure --disable-shared \
 --with-blas="-lmkl_intel_lp64 -lmkl_gnu_thread -lmkl_core -lgomp5 -lpthread -lm -ldl"
\end{verbatim}
The default number of threads will be chosen by the OpenMP software,
but can be controlled by setting \code{OMP\_NUM\_THREADS} or
\code{MKL\_NUM\_THREADS}.

\subsubsection{Using Intel Compilers}

\JAGS\ has been successfully built with the Intel oneAPI compilers. To
set up the environment for using these compilers call the
\file{setvar.sh} shell script, {\em e.g.}
\begin{verbatim}
source /opt/intel/oneapi/setvars.sh
\end{verbatim}
Then call the configure script with the Intel compilers:
\begin{verbatim}
CC=icx CXX=icpx F77=ifx ./configure 
\end{verbatim}

\subsubsection{Using Clang}

\JAGS\ has been built with the clang compiler for C and C++ (version 21.4).
The configuration was
\begin{verbatim}
CC="clang" CXX="clang++ -stdlib=libc++" ./configure
\end{verbatim}
This configuration will use the GNU Fortran compiler \code{gfortran}.
Alternatively, you can the flang compiler by also specifying
\begin{verbatim}
FC=flang F77=flang
\end{verbatim}
on the configure line.

\clearpage
\section{macOS}

A binary distribution of \JAGS\ is provided for Mac OS X versions 
10.9 to 10.11 and macOS 10.12 onwards, which is compatible with the
current binary distribution of \R\ and the corresponding \texttt{rjags} 
and \texttt{runjags} packages that are provided on CRAN. These instructions 
are only for those users that want to install \JAGS\ from source.

The recommended procedure is to build JAGS using clang and the libc++ 
standard library, which have been the default since OS 10.9.  This provides 
compatibility with all builds of \R\ available on CRAN from version 3.3.0 
onwards, as well as ``Mavericks builds'' of earlier versions of \R.  Users 
needing to build against the libstdc++ library and/or with a version of 
Mac OS X predating 10.9 (Mavericks) should refer to the installation 
instructions given in older versions of the \JAGS~manual.

\subsection{Required tools}

If you wish to build from a released source package i.e.
\file{JAGS-\release.tar.gz}, you will need to install command line compilers. 
The easiest way to do this is using the Terminal application from 
\file{/Applications/Utilities} - opening the application gives you a terminal 
with a UNIX shell.  Run the following command on the terminal and follow the
instructions:

\begin{verbatim}
xcode-select --install
\end{verbatim}


You will also need to install the gfortran package, which you 
can download from:

\url{https://gcc.gnu.org/wiki/GFortranBinaries#MacOS}

This setup should be sufficient to build the \JAGS\ sources and also source
packages in R.  All the necessary libraries such as BLAS and LAPACK are
included within macOS.  Additional tools are required to run the optional 
test suite (see section \ref{section:osxtest}).


\subsection{Basic installation}


\subsubsection{Prepare the source code}

Move the downloaded \file{JAGS-\release.tar.gz} package to some suitable
working space on your disk and double click the file.  This will
decompress the package to give a folder called \file{JAGS-\release}.  
You now need to re-open the Terminal and
change the working directory to the \JAGS\ source code. In the Terminal
window after the \$ prompt type \code{cd} followed by a space.  In the Finder
drag the \file{JAGS-\release} folder into the Terminal window and hit return.  If this
worked for you, typing \code{ls} followed by a return will list the contents
of the \JAGS\ folder.


\subsubsection{Set up the environment}
\label{section:osxenvironment}

Before configuring \JAGS\ it is first necessary to set a linker flag to
include the Accelerate framework (\url{https://developer.apple.com/documentation/accelerate}).  
This allows the \JAGS\ installation to use Apple's implementation of BLAS.

Copy and paste the following command into the Terminal window:

\begin{verbatim}
export LDFLAGS="-framework Accelerate"
\end{verbatim}

Other compiler options such as optimisation flags can also be set at 
this stage if desired, for example:

\begin{verbatim}
export CFLAGS="-Os"
export CXXFLAGS="-Os"
export FFLAGS="-Os"
\end{verbatim}

Note that \JAGS\ is usually compiled using \code{-O2} by default, but 
it may be necessary to specify this explicitly depending on the version
of the compiler being used.

\subsubsection{Configuration}

To configure the package type:

\begin{verbatim}
./configure
\end{verbatim}

This instruction should complete without reporting an error.

\subsubsection{Compile}
\label{section:osxcompile}

To compile the code type: 

\begin{verbatim} 
make -j 8 
\end{verbatim} 

The number `8' indicates the number of parallel build threads that
should be used (this will speed up the build process).  In general this
is best as twice the number of CPU cores in the computer - you may want
to change the number in the instruction to match your machine. Again,
this instruction should complete without errors.

\subsubsection{Install}
\label{section:osxinstall}

Finally to install \JAGS\ you need to be using an account with
administrator privileges.  Type: 

\begin{verbatim}
sudo make install
\end{verbatim} 

This will ask for your account password and install the code ready to 
run as described in the User Manual. You need to ensure that
\texttt{/usr/local/bin} is in your PATH in order for the command 
\texttt{jags} to work from a shell prompt.

\subsection{Running the test suite}
\label{section:osxtest}

\subsubsection{Installing CppUnit}

As of \JAGS\ version 4, a test suite is included with the source code that can be
run to ensure that the compiled code produces the expected results.  To run
this code on your installation, you will need to download the CppUnit framework either
using Homebrew (see section \ref{section:osxtools}) or from:

\url{http://freedesktop.org/wiki/Software/cppunit/} 

For the latter, download the source code under ``Release Versions'' corresponding 
to the latest release (currently Cppunit 1.14.0), unarchive the file, and then navigate 
a terminal window to the working directory inside the resulting folder.  Then follow 
the usual terminal commands (as given on the website) to install CppUnit.  

\subsubsection{Running the tests}

The test suite is run following the instructions given in section \ref{section:osxinstall}, 
using the following additional command:

\begin{verbatim} 
make check
\end{verbatim} 

If successful, a summary of the checks will be given.  If compiler errors are encountered, 
you may need to add the following compiler flag (and subsequent reconfiguration) in order to force 
the compiler to build with C++11, as required by CppUnit 1.14.0 and later:

\begin{verbatim} 
export CXXFLAGS="-std=c++11 $CXXFLAGS"
./configure
make check
\end{verbatim}

Note that the configuration step may also need to be repeated if CppUnit 
was not installed the first time this was run.  In this case, you may 
also need to clean the existing compiled code before running \texttt{make check} using:

\begin{verbatim} 
make clean
\end{verbatim} 



\subsection{Tips for developers and advanced users}

\subsubsection{Additional tools}
\label{section:osxtools}

Some additional tools are required to work with code from the \JAGS\ repository.
The easiest way of obtaining the necessary utilities is using Homebrew,
which can be installed by following the instructions at \url{http://brew.sh}

The following instructions have been verified to work with both Mac OS X 10.9 
(Mavericks) and macOS 10.12 (Sierra), and should also work with other
supported versions of OS X / macOS.  If problems are encountered with
these instructions on OS X 10.8 or earlier, then an alternative method
of installing the required tools using e.g. MacPorts (as given in version 4.1.0 of 
the \JAGS\ manual) may be more successful.

\subsubsection{Working with the development code}

If you want to work on code from the \JAGS\ repository, you will need to
build and install the auxillary GNU tools (autoconf, automake and
libtool), as well as mercurial, bison, and flex as follows:

\begin{verbatim}
brew install mercurial
brew install autoconf
brew install automake
brew install libtool
brew install pkg-config
brew install bison
brew install flex
\end{verbatim}

Note that CppUnit can also be installed using the same method:

\begin{verbatim}
brew install cppunit
\end{verbatim}

The following sequence should then retrieve a clone of the current
development branch of \JAGS, and prepare the source code:

\begin{verbatim} 
hg clone -r release-4_patched http://hg.code.sf.net/p/mcmc-jags/code-0
cd code-0
autoreconf -fis
\end{verbatim}

The following modification to the \code{PATH} is also currently 
required to find the Homebrew versions of bison and flex:

\begin{verbatim} 
export PATH="/usr/local/opt/bison/bin:/usr/local/opt/flex/bin:$PATH"
\end{verbatim}

For more information see \code{brew info bison} and \code{brew info flex}

Once these commands have been run successfully, you should be able to follow 
the configuration and installation instructions from section \ref{section:osxenvironment} 
onwards.

\subsubsection{Using ATLAS}

Rather than using the versions of BLAS and LAPACK provided within OS X, it
is possible to use ATLAS, which is available from \url{http://math-atlas.sourceforge.net}
This can be either be installed by following the instructions given at
\url{http://math-atlas.sourceforge.net/atlas_install/}, or by using MacPorts 
(\url{https://www.macports.org/}) with the following Terminal command:

\begin{verbatim}
sudo port install atlas
\end{verbatim}

Once ATLAS is successfully installed, the \texttt{-framework Accelerate} flag can
be omitted from the instructions given in section \ref{section:osxenvironment}.



\clearpage
\section{Windows}
\label{section:windows}

These instructions use MSYS (Minimal System), which is a software
development environment to help build and deploy Unix code on Windows.
You need some familiarity with Unix in order to follow the build
instructions but, once built, \JAGS\ can be installed on any PC
running windows, where it can be run from the Windows command prompt.

\subsection{Preparing the build environment}

You need to install the following packages
\begin{itemize}
\item The \Rtools\ compiler suite for Windows
\item NSIS, including the AccessControl plug-in  
\end{itemize}

\subsubsection{Rtools}

Rtools is a set of compilers and utilities used for compiling R on
Windows. Rtools can be downloaded from your nearest CRAN mirror
(\url{https://cran.r-project.org/bin/windows/Rtools/}). We have to use
the same compiler for R and \JAGS\ to maintain binary compatibility.
For R 4.5.0 and above, the \Rtools\ toolchain is used.

\Rtools\ comes with its own version of MSYS, a minimal unix shell that
we will use to build \JAGS. The MSYS shell does not include the
\Rtools\ compiler tools on the PATH, so this must be added manually
with
\begin{verbatim}
export PATH=/x86_64-w64-mingw32.static.posix/bin:$PATH
\end{verbatim}

\subsubsection{NSIS}

The Nullsoft Scriptable Install System
(\url{https://nsis.sourceforge.io}) allows you to create a
self-extracting executable that installs \JAGS\ on the target PC.
These instructions were tested with NSIS 3.08.  You must also install
the AccessControl plug-in for NSIS, which is available from
\url{http://nsis.sourceforge.io/AccessControl_plug-in}. The plug-in
is distributed as a ZIP file which is unpacked into the installation
directory of NSIS.

At the time of writing, the ZIP file for the AccessControl plug-in is
not correctly configured. Within the Plugins sub-folder of the NSIS
installation directory, you will need to copy
\texttt{i386-ansi/AccessControl.dll} into \texttt{x86-ansi} and
\texttt{i386-unicode/AccessControl.dll} into \texttt{x86-unicode}.

\subsection{Compiling \JAGS}

Unpack the \JAGS\ source
\begin{verbatim}
tar xfvz JAGS-5.0.0.tar.gz
cd JAGS-5.0.0
\end{verbatim}
and configure \JAGS 
\begin{verbatim}
lt_cv_deplibs_check_method=pass_all \
./configure --host=x86_64-w64-mingw32.static.posix 
\end{verbatim}
The \verb+pass_all+ parameter passed to libtool bypasses some file
checks that libtool conducts on dependent libraries. These checks fail
on Windows (presumably due to the different file system) and libtool
falsely concludes that it cannot build a shared library for
JAGS. Bypassing these checks with \verb+pass_all+ ensures that libtool
will build a shared library, as well as dynamic shared objects for the
modules.

After the configure step, type
\begin{verbatim}
make win64-install
\end{verbatim}
This will install \JAGS\ into the subdirectory \file{win/inst64}.  Note
that you must go straight from the configure step to \texttt{make
  win64-install} without the usual step of typing \texttt{make} on its
own.  The \texttt{win64-install} target resets the installation
prefix, and this will cause an error if the source is already
compiled.

You can now create the installer. Make sure that the file
\file{makensis.exe}, provided by NSIS, is in your PATH. For a typical
installation of NSIS, on 64-bit windows:
\begin{verbatim}
PATH=$PATH:/c/Program\ Files\ \(x86\)/NSIS
\end{verbatim}
Then type
\begin{verbatim}
make installer
\end{verbatim}
After the build process finishes, the self extracting archive will be
in the subdirectory \file{win}.

\subsection{Running the unit tests}

In order to run the unit tests on Windows you must first install the
cppunit library from source. Download the file \file{cppunit-1.15.1.tar.gz}
from (\url{http://dev-www.libreoffice.org/src/cppunit-1.15.1.tar.gz}
and unpack it:
\begin{verbatim}
tar xfvz cppunit-1.15.1.tar.gz
cd cppunit-1.15.1
\end{verbatim}
Then compile and install as follows:
\begin{verbatim}
./configure --prefix=/opt --build=x86_64-w64-mingw32.static.posix \
            --disable-shared
make
make install
\end{verbatim}

The configure option \texttt{--disable-shared} prevents the creation
of the DLL \file{libccpunit.dll} and builds only the static library
\file{libcppunit.a}. Without this option, the unit tests will fail.
One of the major limitations of static linking to the C++ runtime is
that you cannot throw exceptions across a DLL boundary.  Linking the
test programs against \file{libcppunit.dll} will result in uncaught
exceptions and apparent failures for some of the tests, so it must be
disabled.

\JAGS\ uses \texttt{pkg-config} to detect \texttt{libcppunit}. This
does not seem to be reliable on Windows but the alternative implementation
\texttt{pkgconf} does work. You need to set the environment variable
\texttt{PKG\_CONFIG}
\begin{verbatim}
export PKG_CONFIG=pkgconf
\end{verbatim}
and then extend the search path
\begin{verbatim}
PKG_CONFIG_PATH=$PKG_CONFIG_PATH:/opt/lib/pkgconfig
\end{verbatim}
before configuring \JAGS.

Then run \code{make check} after \code{make win64-install}.

\subsection{Installing rjags}

To install the \textsf{rjags} package from source you need to let R
know where JAGS is installed. The \texttt{rjags} package will assume
that you have used the default installation directory for \JAGS. If
you have not then you need to set the variable \texttt{JAGS\_ROOT} in
the file \texttt{.R/Makevars.win} in your home directory. Home is
usually the Documents folder on Windows but this can be checked by
calling
\begin{verbatim}
Sys.getenv("HOME")
\end{verbatim}
within an R session. Then edit the file \texttt{.R/Makevars.win} to
give an appropriate value for \texttt{JAGS\_ROOT}, {\em e.g.} 
\begin{verbatim}
JAGS_ROOT=c:\Progra~1\JAGS\JAGS-5.0.0
\end{verbatim}
Note that this is a Make variable, not an environment
variable. Setting \texttt{JAGS\_ROOT} as an environment variable will
have no effect.

\subsection{Running the examples}

The BUGS classic examples (file \file{classic-bugs.tar.gz} from the
JAGS Sourceforge site) can be run from the MSYS shell
using the \command{make} command provided by Rtools. You must have R
installed, along with the packages \code{rjags} and \code{coda}, both
of which are available from CRAN (\url{cran.r-project.org}).

It is necessary to add R to the search path.
\begin{verbatim}
PATH=$PATH:/c/Program\ Files/R/R-4.5.2/bin
\end{verbatim}
Then
\begin{verbatim}
tar xfvz classic-bugs.tar.gz
cd classic-bugs
cd vol1
make Rcheck
\end{verbatim}
will check all examples in volume 1 using the \code{rjags} package. Repeat
for \file{vol2} to complete the checks.

\end{document}
